% Methodology (P8.6)
\section{Methodology}\label{sec:methodology}

The analysis runs as one pipeline (\texttt{python -m src.pipeline}). Every stage reads and
writes files with a fixed schema, and every parameter lives in \texttt{config/config.yaml}.

\subsection{Analysis grid and features}

The study area plus a \SI{3}{\km} buffer is divided into \SI{100}{\m} cells aligned with the
\SI{10}{\m} raster grid, so every cell holds exactly $10\times10$ land-cover pixels. The
58{,}218 cells whose centres lie inside the study area are analysed. For each cell and year
the feature table holds:
\begin{itemize}
  \item \textbf{own-cell land cover}: the shares of water, trees, flooded vegetation, crops,
        built area, bare ground and rangeland among its 100 pixels;
  \item \textbf{context rings}: the same shares in square rings of \SI{250}{\m} and
        \SI{500}{\m} around the cell, \emph{excluding} the cell itself;
  \item \textbf{distances} from the cell centre to the nearest built-up pixel outside the cell,
        to water, to a major road and to any road (log-transformed, capped at \SI{3}{\km}, the
        buffer width, so that edge cells see everything within the cap);
  \item \textbf{road density}: road length within \SI{500}{\m} per km\textsuperscript{2};
  \item \textbf{terrain}: mean elevation, mean and maximum slope;
  \item building count and footprint share (kept for analysis only).
\end{itemize}

\textbf{Leakage rule.} A model trained to recognise built land could simply read it from the
cell's own land cover, whose shares sum to one minus the built share, or from its buildings.
In a prototype, own-cell shares alone separated built from non-built cells perfectly. The
similarity and the random forests therefore use only the context features: rings, distances,
road density and terrain (20 inputs). The code refuses any other column.

\subsection{Land classification}

K-Means \citep{pedregosa2011} groups the cells on seven standardised inputs: the shares of water,
trees, crops, built area and rangeland, mean elevation and mean slope. (Snow is absent and
flooded vegetation and bare ground are almost always zero; maximum slope correlates at
$r=0.95$ with mean slope.) The number of clusters is chosen from $k=3$ to 10 with inertia, the
silhouette \citep{rousseeuw1987} and the Davies--Bouldin index \citep{davies1979}; Gaussian
mixtures are compared. Clusters become classes by rules on their centroids:
water-dominated $\rightarrow$ excluded; built-up if built land exceeds 30\,\% and is the largest
land group; forest if trees exceed 50\,\%, or trees and rangeland together exceed 50\,\% on
slopes above \ang{8} (the Bannerghatta scrub forest, which the ESRI map partly calls rangeland);
otherwise usable, with cropland flagged as agricultural.

The \textbf{exclusion mask} is applied last, cell by cell: more than half of the cell is water,
flooded vegetation or snow in either year of the run, or an OSM lake; mean slope above
\ang{15}; more than half no-data; or more than half inside a protected area.

\subsection{Building similarity}\label{sec:similarity}

\paragraph{Reference buildings.} Cells with an OpenStreetMap building, plus ESRI built cells
where OSM has none: OSM maps 62--75\,\% of built cells in the north but about 1\,\% in the south
in 2018, so OSM alone would describe only northern surroundings.

\paragraph{Profiles.} A location's profile is its cell's 20 context features, standardised with
the mean and standard deviation of the reference and candidate cells of the run.

\paragraph{Single query.} For building A, every usable cell is ranked by the Euclidean distance
between its profile and A's (cosine distance is available for comparison). Each result lists
the features where B is closest to A and the ones where it differs most, in standard deviations.

\paragraph{Aggregate similarity.} For every candidate cell, the mean distance to its ten nearest
reference profiles, never counting the cell itself, is rescaled to 0--1 between the 1st and
99th percentile (higher = more similar).

\subsection{Learned models}

Two random forests \citep{breiman2001} use the 20 context features of 2018:
\begin{itemize}
  \item the \textbf{change RF} learns from the cells that were non-built in 2018 and 2019 and
        became built in 2020 and 2021 (355 positives among 15{,}658 candidates);
  \item the \textbf{state RF} learns built vs non-built cells in 2018 (for comparison).
\end{itemize}
Both are class-balanced (300 trees). Hyperparameters are tuned by spatial cross-validation
\citep{roberts2017, valavi2019}: cells are grouped into \SI{2}{\km} blocks, which are dealt into
five folds, so no block is in both training and test data. Tuning uses only the training label.
Features enter as percentiles within their year, because OSM road mapping grew a lot between
the 2018 and current snapshots; within a year the transform keeps each tree's partition of its
training data. Importance is reported both by impurity and by permutation on the held-out folds.

\subsection{MCDA / AHP baseline}

The classic overlay rescales five criteria to 0--1: distance to built-up land (1 next to it, 0
beyond \SI{2}{\km}), distance to a major road (0 beyond \SI{3}{\km}), slope (0 at \ang{15}),
own-cell land cover (rangeland and bare ground 1, crops 0.6, trees 0.2) and a water setback (0
within \SI{100}{\m}, 1 beyond \SI{200}{\m}). Weights come from an AHP pairwise matrix
\citep{saaty1980}: built-up and road distance 0.30 each, slope and land cover 0.16 each, water
0.09, with a consistency ratio of 0.003. MCDA is not trained, so its use of own-cell land cover
is not the leakage problem above.

\subsection{Combined score, classes and candidate sites}

The final score is a weighted mean of the aggregate similarity and the change RF. Each is first
converted to its percentile among the scored cells, because the two have very different
distributions. The weight is chosen on the validation run, but on the change RF's training label
(growth by 2020/21), never on the later growth used for validation, and then frozen: 0.3 for the
similarity and 0.7 for the change RF.

Usable cells are classed by quantiles of the score: S1 at least the 90th percentile, S2 the 70th,
S3 the 40th, otherwise N. Forest, built-up and excluded land is always N. Edge-adjacent S1 cells
are merged into candidate patches of at least \SI{1}{\hectare}, each with its area, mean score,
slope, distances and land cover. The top sites are the 20 best patches of at least
\SI{5}{\hectare}.

\subsection{Implementation}

Everything is written in Python: rasterio and GeoPandas for the data, scikit-learn for the
models, Matplotlib and Folium for the maps; OpenStreetMap data comes through OSMnx
\citep{boeing2017}. Each stage has unit tests on synthetic data, a phase gate on the real data,
and an independent check script that recomputes its output another way (for example pixel
values against the source rasters, brute-force distances, or AUC from ranks). The whole pipeline
runs in about 12 minutes after the first download, and a fresh clone reproduces the validation
run exactly.

\paragraph{Live demo.} The results are also published as a static website:
\begin{center}
  \url{https://iamabhinav-dev.github.io/gis-urban-suitability/}
\end{center}
It shows the suitability map with the top 20 sites, an interactive map of every layer, and a
``find places like A'' tool. The tool runs the single-building query of Section~\ref{sec:similarity}
in the browser on the exported cell profiles; a check script confirms that it returns the same
top 20 locations as the Python query.
